\section{Seção com Símbolos, Notação e Tabela}

% Gera uma caixa de sumário local das subseções desta seção
\gerarSumarioLocal

Texto com símbolos: $\forall x \in \mathbb{R}$, $\exists y>0$ tal que $|x|<y$ e
$\sum_{k=0}^{n} k = \frac{n(n+1)}{2}$.

\[
\lim_{n\to\infty}\left(1+\frac{1}{n}\right)^n = e.
\]

\subsection{Tabela de Simulação}

% \tabelaSimulacao{<caption>}{<label>}{<linhas>}
\tabelaSimulacao
{Cenário 1 -- Acerto das 16 jogadas (Simulação)}
{tab:cenario1}
{
    1  & Resetar circuito & \texttt{reset\_in = 1} & FSM no estado inicial & \\
    2  & Aguardar estabilidade & - & - & \\
    3  & Acionar iniciar & \texttt{iniciar\_in = 1} & Aguardando jogada 1 & Aguardando jogada 1 \\
    % ... adicione as outras linhas aqui ...
    19 & Acionar jogada 16 & \texttt{chaves = 4'b0100} &
    \begin{tabular}[c]{@{}l@{}}Sinais \textbf{acertou} e \\ \textbf{pronto} ativados\end{tabular} & \\
}

\subsection{Tabela de Duas Colunas}

% \tabeladuas{<caption>}{<label>}{<col1>}{<col2>}{<linhas>}
\tabeladuas
  {Resumo de parâmetros do experimento}
  {tab:parametros}
  {Parâmetro}
  {Valor}{
    Tamanho de palavra & 16 bits \\
    Clock & 50\,MHz \\
    Nº de testes & 16 \\
    Ambiente & Simulação (Icarus/GTKWave) \\
  }

\subsection{Tabela de Três Colunas}

% \tabelatres{<caption>}{<label>}{<col1>}{<col2>}{<col3>}{<linhas>}
\tabelatres
  {Resumo do experimento}
  {tab:resumo-exp}
  {Item}
  {Descrição}
  {Observação}{
    Clock & 50\,MHz & FPGA DE10-Lite \\
    Entradas & 4 bits (\texttt{chaves}) & Usuário \\
    Saídas & \texttt{acertou}, \texttt{errou} & LEDs \\
  }
